% =====================================================================
%  Telas do etilômetro DAD01 desenhadas em vetor (TikZ)
%  Baseadas no firmware v6.4.8: tela 320 x 480 px (retrato), textos de
%  device/lang (PT), posições e cores de objects/display e objects/test.
%  Coordenadas em pixels da tela, y para baixo. Para legibilidade no papel
%  as fontes foram ampliadas e textos longos quebrados em mais linhas.
% =====================================================================

% Cores do display (RGB565 do firmware convertido para RGB)
\definecolor{tftgreen}{RGB}{22,178,52}    % TFT_GREEN (tom do LCD)
\definecolor{tftred}{RGB}{214,24,24}      % TFT_RED
\definecolor{tftorange}{RGB}{240,160,0}   % TFT_ORANGE
\definecolor{tftyellow}{RGB}{33,157,90}   % YELLOW  (0x24EB)
\definecolor{tftdgrey}{RGB}{38,37,38}     % DARK_GREY (0x2104)
% Moldura (cor da carcaça)
\definecolor{dadbody}{HTML}{3AA64C}
\definecolor{dadbodyedge}{HTML}{1A9530}
\definecolor{dadlogo}{HTML}{318E41}

% Fontes Roboto (as mesmas do display)
\newcommand{\fA}{\fontsize{12}{13}\roboto\selectfont}                      % títulos
\newcommand{\fB}{\fontsize{11.5}{12.5}\roboto\bfseries\selectfont}         % botões
\newcommand{\fC}{\fontsize{10}{11}\roboto\robotocondensed}                 % textos auxiliares

\tikzset{dadpx/.style={x=0.125mm,y=-0.125mm,every node/.style={inner sep=0pt,outer sep=0pt,anchor=center,align=center}}}

% \dad[largura]{conteúdo da tela}
\newcommand{\dad}[2][27mm]{%
  \resizebox{#1}{!}{%
  \begin{tikzpicture}[dadpx]
    % bocal lateral
    \draw[fill=dadbody,draw=black!75,line width=0.5pt,rounded corners=0.8mm] (-62,44) rectangle (-40,114);
    % corpo
    \draw[fill=dadbody,draw=black!75,line width=0.6pt,rounded corners=4mm] (-40,-50) rectangle (360,690);
    % moldura da tela
    \draw[fill=black!85,draw=black!75,line width=0.4pt,rounded corners=1.2mm] (-12,-14) rectangle (332,494);
    \begin{scope}
      \clip (0,0) rectangle (320,480);
      \fill[black] (0,0) rectangle (320,480);
      #2
    \end{scope}
    % logotipo em baixo relevo
    \node[text=dadlogo,font=\fontsize{30}{30}\roboto\bfseries\itshape\selectfont] at (160,592) {SIGHIR};
  \end{tikzpicture}}}

% Barra de progresso (LoadingWindow::drawBar): y = topo, #2 = fração 0..1
\newcommand{\dadbarra}[2]{%
  \fill[white] (32,#1) rectangle (288,{#1+26});
  \fill[black] (35,{#1+3}) rectangle (285,{#1+23});
  \fill[white] (35,{#1+3}) rectangle ({35+250*#2},{#1+23});}

% Aviso legal exibido durante câmera e análise (TestScreens::inmetro)
\newcommand{\dadinmetro}{%
  \node[white,font=\fC] at (160,178) {Este dispositivo não se};
  \node[white,font=\fC] at (160,212) {enquadra no âmbito};
  \node[white,font=\fC] at (160,246) {probatório de fiscalização};
  \node[white,font=\fC] at (160,280) {ou qualquer outro tipo};
  \node[white,font=\fC] at (160,314) {de punição legal};
  \node[white,font=\fC] at (160,410) {Portaria INMETRO Nº369};}

% ---------------- Telas ----------------
\newcommand{\telaIniciado}{%
  \node[white,font=\fA] at (160,225) {Teste Iniciado};
  \node[white,font=\fC] at (160,322) {clique na tela para};
  \node[tftred,font=\fC] at (160,352) {Adiamento de Emergência};
  \node[white,font=\fC] at (160,433) {iniciando em: 3};}

\newcommand{\telaCamera}{%
  \node[white,font=\fA] at (160,44) {Aguardando Câmera};
  \dadbarra{76}{0.4}
  \node[white,font=\fC] at (160,128) {40\%};
  \dadinmetro}

\newcommand{\telaCalibrando}{%
  \node[tftred,font=\fC] at (160,46) {21480 - 27 ºC};
  \node[white,font=\fA] at (160,190) {Calibrando};
  \node[white,font=\fA] at (160,226) {Sensor};
  \dadbarra{264}{0.68}
  \node[white,font=\fC] at (160,316) {68\%};
  \node[white,font=\fC] at (160,380) {Toque Para Adiar};}

\newcommand{\telaAssopre}{%
  \node[white,font=\fontsize{15}{16}\roboto\selectfont] at (160,225) {Assopre};}

\newcommand{\telaAnalisando}{%
  \node[white,font=\fA] at (160,40) {Analisando};
  \dadbarra{70}{0.55}
  \node[white,font=\fC] at (160,122) {55\%};
  \dadinmetro}

\newcommand{\telaResultado}{%
  \node[white,font=\fontsize{14}{15}\roboto\selectfont] at (160,225) {0.000 mg/L};
  \node[tftred,font=\fA] at (160,423) {21875};}

\newcommand{\telaLiberado}{%
  \fill[tftgreen] (0,0) rectangle (320,480);
  \node[white,font=\fB] at (160,205) {Veículo};
  \node[white,font=\fB] at (160,245) {Desbloqueado};}

\newcommand{\telaAlcool}{%
  \fill[tftred] (0,0) rectangle (320,480);
  \node[white,font=\fB] at (160,205) {Álcool};
  \node[white,font=\fB] at (160,245) {Detectado};}

\newcommand{\telaOpcoes}{%
  \fill[tftyellow] (0,0) rectangle (320,160);
  \fill[tftdgrey] (0,160) rectangle (320,320);
  \fill[white] (0,320) rectangle (320,480);
  \node[white,font=\fB] at (160,80) {Contrassenha};
  \node[tftyellow!70!white,font=\fB] at (160,222) {Realizar};
  \node[tftyellow!70!white,font=\fB] at (160,258) {Outro Teste};
  \node[black,font=\fB] at (160,382) {Desligar};
  \node[black,font=\fB] at (160,418) {Display};}

\newcommand{\dadtecla}[3]{% x, y, rótulo
  \fill[white] (#1,#2) rectangle ({#1+50},{#2+50});
  \node[black,font=\fB] at ({#1+25},{#2+25}) {#3};}

\newcommand{\telaSenha}{%
  \node[white,font=\fB] at (165,40) {Contrassenha};
  \fill[white] (40,60) rectangle (290,105);
  \node[black,font=\fB] at (165,83) {4721};
  \node[white,font=\fB] at (165,130) {Digite};
  \fill[white] (40,150) rectangle (290,195);
  \dadtecla{80}{220}{1}\dadtecla{140}{220}{2}\dadtecla{200}{220}{3}
  \dadtecla{80}{280}{4}\dadtecla{140}{280}{5}\dadtecla{200}{280}{6}
  \dadtecla{80}{340}{7}\dadtecla{140}{340}{8}\dadtecla{200}{340}{9}
  \dadtecla{80}{400}{<}\dadtecla{140}{400}{0}\dadtecla{200}{400}{+}
  \dadtecla{260}{340}{\^{}}\dadtecla{260}{400}{OK}}

\newcommand{\telaRandomico}{%
  \fill[tftyellow] (0,0) rectangle (320,240);
  \fill[tftdgrey] (0,240) rectangle (320,480);
  \node[white,font=\fB] at (160,120) {Iniciar Teste};
  \node[tftyellow!70!white,font=\fB] at (160,360) {Adiar};}

\newcommand{\telaMotorista}{%
  \fill[tftred] (0,0) rectangle (320,480);
  \node[white,font=\fB] at (160,190) {MOTORISTA};
  \node[white,font=\fB] at (160,226) {NÃO};
  \node[white,font=\fB] at (160,262) {AUTORIZADO};
  \node[white,font=\fC] at (160,360) {Estacione o veículo};}

% Tela com legenda numerada embaixo: \dadfig[largura]{tela}{nº}{legenda}
\newcommand{\dadfig}[4][27mm]{%
  \begin{minipage}[t]{#1}\centering
    \dad[#1]{#2}\\[1.5mm]
    {\footnotesize\textbf{\color{sgreendark}#3}\enspace\color{sdark}#4}
  \end{minipage}}
